
\begin{proof}
By definition,
\[
\langle \Theta O,O\rangle_{\mathrm{rel}}
=
\frac{1}{Z_{\mathrm{ref}}[0,0,0]}\,
\langle \Theta O,O\rangle_{\mathrm{adm}}.
\]
The denominator is a fixed positive normalizing constant, so the sign is determined entirely by the
admissible numerator.

Split $O$ into its bosonic and fermionic parts on the retained admissible sector. For the bosonic
part, the Bernstein representation of the admissible kernel writes the Euclidean covariance as a
positive superposition of reflection-stable heat kernels. Reflection across $\tau=0$ therefore sends
every positive-time bosonic observable to a positive quadratic form, so the bosonic contribution to
\[
\langle \Theta O,O\rangle_{\mathrm{adm}}
\]
is nonnegative.

For the fermionic part, the CAR reflection-positivity hypothesis on the admissible subspace implies
that the reflected two-point kernel is positive in the standard fermionic sense; equivalently, the
corresponding Pfaffian / determinant form on positive-time test spinors is nonnegative. Because the
bosonic and fermionic admissible factors commute inside the normalized expectation, their product
remains nonnegative. Dividing by the fixed reference normalization preserves that sign. Hence
\[
\langle \Theta O,O\rangle_{\mathrm{rel}}\ge 0.
\]
\end{proof}

\begin{theorem}[Operative admissibility selector]
Under \TFPTxref{def:full-admissibility-complex,thm:padm-upgrade,prop:master-closure-roadmap}, the
operator $\Padm$ is idempotent on the retained branch and realizes the admissibility predicate
on all sectors used by the closed output statements:
\[
\mathsf{Adm}(X)=1
\quad\Longleftrightarrow\quad
\Padm X = X.
\]
\end{theorem}

\begin{proof}[Proof sketch]
By \TFPTxref{thm:padm-upgrade}, $\Padm=\Pprim P_{\mathrm{sing}}P_\Theta$ is the orthogonal projector
onto the harmonic admissible sector of the full complex and is therefore idempotent by
construction. \TFPTxref{cor:factorized-admissibility-selector} recovers the explicit factorized
selector. The predicate and operator languages agree once the positive transport sector is fixed.
\end{proof}

\begin{theorem}[Topological family closure on the admissible branch]
By \TFPTxref{thm:master-closure,thm:topological-family-closure,cor:boundary-bulk-occupancy},
\[
\chi(X_f^\circ)=-2,
\qquad
F\cong\mathbb C^3,
\qquad
N_{\mathrm{fam}}=3.
\]
Consequently
\[
\Oadm=N_{\mathrm{fam}}\dim S^+=3\times 16=48,
\qquad
\deltatop=\Oadm\,\cthree^4.
\]
\end{theorem}

\begin{proof}[Proof sketch]
This theorem is now a direct consequence of the master closure theorem from Section~6. The
topological family closure and the orbifold corner count have already been compressed there
into one statement, so the present layer merely records the corresponding admissibility
consequence.
\end{proof}

\begin{theorem}[Bosonic index and local spectral-scale closure]
On the admissible branch,
\[
\Ind_{\mathrm{rel}}^{\Padm}(\mathcal B_{E_2}^+)=1,
\qquad
\Ind_{\mathrm{rel}}^{\Padm}(\mathcal B_{E_3}^+)=0,
\]
so the unique light seam-even bosonic zero mode is
\[
\Phi\in\Gamma(E_2)\cong(1,2)_{1/2},
\qquad
\taudbl^{\ast}\Phi=+\Phi.
\]
Moreover the same closure layer yields
\[
\begin{aligned}
\frac{\Mpl^2}{\lambda_\Sigma^2}&=\frac{\rho_\star}{2\pi^2},
&
G_N\lambda_\Sigma^2&=\frac{\pi}{4\rho_\star},
\\
\frac{\vgeo}{\Mpl}&=\gcar\,\beta_{\mathrm{rad}}^2
\exp\!\left[-\frac{\alpha^{-1}(0)+\delta_{\mathrm{ph}}}{5}\right],
\\
G_N\vgeo^2&=\frac{1}{8\pi}\gcar^2\beta_{\mathrm{rad}}^4
\exp\!\left[-\frac{2(\alpha^{-1}(0)+\delta_{\mathrm{ph}})}{5}\right].
\end{aligned}
\]
Here $\vgeo$ is the seam-even UV source scale. The physical electroweak benchmark is the residue
matched quantity $v_{\mathrm{phys}}$ of \TFPTxref{thm:ew-geometric-to-physical-matching}.
\end{theorem}

\begin{proof}[Proof sketch]
The bosonic sector is treated as a genuine Callias / Dirac / Dolbeault-type index problem rather
than a scalar counting exercise. Relative heat-kernel closure together with
\TFPTxref{prop:canonical-bosonic-normalization,lem:two-sheet-zero-mode-normalization,thm:absolute-spectral-planck-closure}
produces the boundary-normalized Einstein branch, the canonical Higgs normalization, and the
seam-even zero-mode formula for $\vgeo$.
\end{proof}

\begin{definition}[Primitive Yukawa generator]
\label{def:primitive-yukawa-generator}
Let
\[
\mathcal I_Y
:=
\bigoplus_{p,q,r,t}
\operatorname{Hom}_{G_{\mathrm{car}}}
\Bigl(
(\Lambda^p E_-\otimes \Lambda^q E_+)
\otimes
(\Lambda^r E_-\otimes \Lambda^t E_+)
\otimes E_+,\,
\mathbb C
\Bigr)
\]
be the graded algebra of carrier-invariant Yukawa trilinears. A homogeneous nonzero invariant is
called primitive if it is indecomposable, equivalently if it does not lie in the ideal generated by
positive-degree invariants of strictly lower carrier degree.
\end{definition}

\begin{lemma}[Lowest primitive Yukawa generator on the closed branch]
\label{lem:lowest-primitive-yukawa-generator}
Assume $\dim E_+=2$ and let $\mathbb Y_{\mathrm{br}}$ be the branch Yukawa tensor of
\TFPTxref{thm:exact-transport-closure,thm:rigid-family-local-system}. Then the primitive generator of
$\mathcal I_Y$ with two nontrivial fermionic legs is unique up to exchange of the two fermionic
legs and belongs to
\[
\operatorname{Hom}_{G_{\mathrm{car}}}
\bigl((E_-\otimes E_+)\otimes \Lambda^2 E_-\otimes E_+,\,\mathbb C\bigr).
\]
\end{lemma}

\begin{proof}
The positive-block determinant neutrality gives $q+t+1=2$, hence $q+t=1$. Up to exchanging the two
fermionic legs, $(q,t)=(1,0)$. Evenness of both fermionic legs then forces $p$ odd and $r$ even.
If $p\ge 3$, the corresponding invariant factors through an extra negative spectator wedge and
therefore lies in the ideal generated by lower carrier degree. Likewise, if $r\ge 4$, the
invariant factors through a spectator bivector insertion and is again decomposable. Primitive
indecomposability therefore forces $p=1$ and $r=2$.
\end{proof}

\begin{theorem}[Carrier rigidity from the primitive component of the branch Yukawa tensor]
\label{thm:yukawa-forces-32}
Let
\[
E=E_-\oplus E_+
\]
be the two-factor carrier and let
\[
G_{\mathrm{car}}:=S\bigl(U(E_-)\times U(E_+)\bigr),
\qquad
S^+=\Lambda^{\mathrm{even}}E.
\]
Let $\mathbb Y_{\mathrm{br}}$ denote the branch Yukawa tensor induced by the canonical transport
kernel and rigid family holonomy of
\TFPTxref{thm:exact-transport-closure,thm:rigid-family-local-system}. Decompose its carrier-homogeneous
components as trilinear maps of the form
\[
B_{p,q;r,t}:
(\Lambda^p E_-\otimes \Lambda^q E_+)
\otimes
(\Lambda^r E_-\otimes \Lambda^t E_+)
\otimes E_+
\longrightarrow
\mathbb C.
\]
Let the primitive component of $\mathbb Y_{\mathrm{br}}$ be the unique indecomposable generator of
\Cref{def:primitive-yukawa-generator,lem:lowest-primitive-yukawa-generator} with two nontrivial
fermionic legs. Then, up to exchanging the two fermionic legs, that primitive component is
necessarily of type
\[
(E_-\otimes E_+) \otimes \Lambda^2 E_- \otimes E_+
\longrightarrow
\Lambda^3 E_-\otimes \Lambda^2 E_+\cong\mathbb C,
\]
and therefore
\[
\dim E_+=2,
\qquad
\dim E_-=3.
\]
Consequently
\[
\dim \Lambda^2 E_+=1,
\qquad
\Lambda^2 E_- \cong E_-^\vee,
\qquad
\dim \Lambda^{\mathrm{even}}E=16.
\]
Hence the rigid carrier is
\[
E=E_3\oplus E_2,
\qquad
X=-2P_3+3P_2,
\qquad
Y=-\frac13 P_3+\frac12 P_2.
\]
\end{theorem}

\begin{proof}
By \TFPTxref{cor:bosonic-rank-two}, the compact Higgs index has already fixed
\[
\dim E_+=2.
\]
Write the primitive nonzero Yukawa component with two nontrivial fermionic legs as
\[
B_{p,q;r,t}:
(\Lambda^p E_-\otimes \Lambda^q E_+)
\otimes
(\Lambda^r E_-\otimes \Lambda^t E_+)
\otimes E_+
\longrightarrow
\mathbb C
\]
and note that \Cref{lem:lowest-primitive-yukawa-generator} forces primitiveness in the invariant-ring
sense to have
\[
(q,t)=(1,0),
\qquad
(p,r)=(1,2)
\]
up to exchanging the two fermionic legs. Thus the primitive carrier component is exactly
\[
(E_-\otimes E_+) \otimes \Lambda^2 E_- \otimes E_+.
\]
Its negative factor contracts canonically as
\[
E_- \otimes \Lambda^2 E_- \longrightarrow \Lambda^3 E_-,
\]
so scalarity forces $\Lambda^3 E_-$ to be the top exterior power. Hence
\[
\dim E_-=3.
\]
With
\[
(\dim E_-,\dim E_+)=(3,2),
\]
one obtains the carrier-signature consequences
\[
\Lambda^2 E_- \cong E_-^\vee,
\qquad
\dim\Lambda^2 E_+=1,
\]
and therefore
\[
\dim \Lambda^{\mathrm{even}}(E_3\oplus E_2)=2^{5-1}=16.
\]
Finally the unimodular primitive two-point generator is fixed by
\[
3q_-+2q_+=0,
\qquad
q_-<0<q_+,
\qquad
\gcd(q_-,q_+)=1,
\]
so $(q_-,q_+) = (-2,3)$ and hence
\[
X=-2P_3+3P_2,
\qquad
Y=\frac{X}{6}=-\frac13 P_3+\frac12 P_2.
\]
\end{proof}

\begin{lemma}[Uniqueness of the cubic weak invariant]
\label{lem:unique-cubic-weak-invariant}
Let $E=E_-\oplus E_+$ be a two-factor carrier on the retained branch and let
\[
\Phi\in E_+
\]
be the unique seam-even light boson selected by the compact bosonic index. Then
\[
\dim \operatorname{Hom}_{G_{\mathrm{car}}}\bigl((E_-\otimes E_+)\otimes \Lambda^2 E_-\otimes E_+,\,\mathbb C\bigr)=1
\]
if and only if
\[
(\dim E_-,\dim E_+)=(3,2).
\]
In particular, on the retained branch the up-type cubic invariant is not an extra choice but the
unique admissible cubic invariant.
\end{lemma}

\begin{proof}
By Schur functor decomposition,
\[
\begin{aligned}
\operatorname{Hom}_{G_{\mathrm{car}}}\bigl((E_-\otimes E_+)\otimes \Lambda^2 E_-\otimes E_+,\mathbb C\bigr)
&\cong
\operatorname{Hom}_{SU(E_-)}(E_-\otimes \Lambda^2E_-,\mathbb C)
\\
&\qquad\otimes
\operatorname{Hom}_{SU(E_+)}(E_+\otimes E_+,\mathbb C).
\end{aligned}
\]
The first factor is nonzero iff $\Lambda^2E_-\cong E_-^\vee$, hence iff $\dim E_-=3$. The second
factor is nonzero iff the alternating contraction on two fundamentals is one-dimensional, hence iff
$\dim E_+=2$. In that case both factors are one-dimensional, so the total invariant space is
one-dimensional.
\end{proof}

\begin{corollary}[Carrier-signature premises are discharged internally]
\label{cor:yukawa-discharges-kpp}
Combine \TFPTxref{cor:bosonic-rank-two,thm:yukawa-forces-32,lem:unique-cubic-weak-invariant} with the
one-family decomposition of
$S^+=\Lambda^{\mathrm{even}}(E_3\oplus E_2)$ and with the unique seam-even Higgs theorem. Then the
weak antisymmetric line, the color-dual antisymmetric sector, and the retained light-packet
identification are no longer primitive inputs once the uniqueness of the retained cubic invariant
has been proved on the same branch.
\end{corollary}

\begin{figure}[t]
\centering
\scriptsize
\resizebox{\linewidth}{!}{%
\begin{tikzpicture}[
    >=Latex,
    source/.style={draw=blue!60!black, fill=blue!6, rounded corners, text width=2.75cm, minimum height=0.95cm, align=center},
    force/.style={draw=orange!75!black, fill=orange!8, rounded corners, text width=2.9cm, minimum height=0.95cm, align=center},
    result/.style={draw=green!50!black, fill=green!8, rounded corners, text width=2.9cm, minimum height=0.95cm, align=center},
    note/.style={font=\footnotesize\itshape, text=black!60, align=center}
]
\node[source] (fam) at (-5.2,0.95) {$X_f^\circ,\ N_{\mathrm{fam}}=3$\\family closure};
\node[source] (higgs) at (-1.7,0.95) {$\Ind(\mathcal B_{E_2}^+)=1$\\$\dim E_+=2$};
\node[force] (yuk) at (1.8,0.95) {primitive indecomposable\\Yukawa generator};
\node[result] (car) at (5.3,0.95) {$E_3\oplus E_2,\ g=5,$\\$S^+=\Lambda^{\mathrm{even}}E$};

\node[result] (cons) at (1.8,-1.2) {derived check\\$\Oadm=48=4\cdot 12$};

\draw[->, thick] (fam) -- (higgs);
\draw[->, thick] (higgs) -- (yuk);
\draw[->, thick] (yuk) -- (car);
\draw[->, thick] (car) -- (cons);
\draw[->, thick, dashed] (fam) -- (cons);

\node[note] at (0,-2.35) {Carrier closure runs through Higgs rank and branch Yukawa rigidity; the
old occupancy formula is only a downstream consistency identity.};
\end{tikzpicture}%
}
\caption{Carrier closure after the proof-tree cut. The retained branch first fixes $\dim E_+=2$
through the compact Higgs index, then $\dim E_-=3$ through the primitive indecomposable generator of
the actual branch Yukawa tensor; only afterwards is $\Oadm=48=4\cdot12$ read as a consistency check.}
\label{fig:carrier-closure-cut}
\end{figure}

\begin{corollary}[Hard holonomy form of the Yukawa matrices]
On the admissible branch, each fermion sector obeys the canonical factorization
\[
Y_f=D_{L,f}U_f^\star D_{R,f},
\qquad
U_f^\star\in SU(3)_F,
\]
and, in the family basis fixed by \TFPTxref{thm:exact-transport-closure},
\[
(Y_f)_{ij}
=
\lambdaY^{L^L_{f,i}+L^R_{f,j}}
\left\langle e_i,\mathrm{Hol}_{\Gamma_{ij}^{\min}}(\nabla_F^\star)e_j\right\rangle.
\]
Any factorization by $U_f^\star$ is therefore a consequence of the hard holonomy closure, not an
independent flavor input. The diagonal path-length rows \TFPTxref{eq:path-length-yukawa} are the
singular-value reduction of this matrix statement.
\end{corollary}

\begin{proof}[Proof sketch]
The finite $C_6$ transport geometry fixes the graph-distance weights, while the rigid
$D_4$-equivariant family connection fixes the holonomy factors on the canonical path frame. The
diagonal ladder formulas arise once left and right admissible words are paired on the same
singular-value channel, so the matrix factorization is recovered as a corollary of one closed
holonomy statement rather than as a separate ansatz.
\end{proof}

\begin{definition}[Center-neutral admissible algebra]
Let
\[
\Asing
:=
\overline{\mathrm{span}}\{A_w:\ z(w)=1,\ \varepsilon(w)=+1,\ \Padm A_w=A_w\}.
\]
\end{definition}

\begin{theorem}[Hadronic admissibility closure]
On the admissible branch, the physical hadronic Hilbert space is
\[
\Hhad:=\overline{\Asing|\Omega\rangle}.
\]
All physical hadronic states lie in $\Hhad$, while color-nonsinglet free-word sectors are removed by $\Padm$.
\end{theorem}
